\documentclass[12pt]{article}
\usepackage{amsmath, amssymb}
\usepackage{geometry}
\geometry{margin=1in}

\title{Lecture-22 Scribe: Inequalities and Tail Bounds (From Markov’s Inequality Onward)}
\date{}

\begin{document}

\maketitle

\textit{(Exam Preparation Notes – Step-by-Step Reasoning)}

\hrule

\section*{1. Markov’s Inequality}

\subsection*{1.1 Definition}

Let $X$ be a non-negative random variable $(X \ge 0)$.

For any $a > 0$,

\[
P (X \ge a) \le \frac{E[X]}{a}
\]

\subsection*{1.2 Step-by-Step Proof}

\textbf{Step 1: Start from expectation}

\[
E[X] = \sum_x x \cdot P (X = x)
\]

\textbf{Step 2: Split based on threshold $a$}

\[
E[X] = \sum_{x<a} xP (X = x) + \sum_{x \ge a} xP (X = x)
\]

\textbf{Step 3: Bound the large-value part}

For all $x \ge a$,

\[
x \ge a \Rightarrow xP (X = x) \ge aP (X = x)
\]

Thus,

\[
\sum_{x \ge a} xP (X = x) \ge a \sum_{x \ge a} P (X = x)
\]

\[
= a \cdot P (X \ge a)
\]

\textbf{Step 4: Combine}

\[
E[X] \ge a \cdot P (X \ge a)
\]

\textbf{Step 5: Rearranging}

\[
P (X \ge a) \le \frac{E[X]}{a}
\]

\subsection*{1.3 Example}

If:

\[
E[X] = 10
\]

Then:

\[
P (X \ge 50) \le \frac{10}{50} = 0.2
\]

\subsection*{1.4 Key Insight}

No assumption on distribution.\\
Only requires non-negativity.\\
Provides a loose but universal bound.

\section*{2. Chebyshev’s Inequality}

\subsection*{2.1 Motivation}

Markov’s inequality only uses expectation.

To get a tighter bound, we incorporate variance.

\subsection*{2.2 Definition}

For any random variable $X$ with mean $\mu$ and variance $\sigma^2$,

\[
P (|X - \mu| \ge a) \le \frac{\sigma^2}{a^2}
\]

\subsection*{2.3 Proof Using Markov’s Inequality}

\textbf{Step 1: Define new variable}

Let:

\[
Y = (X - \mu)^2
\]

$Y \ge 0$, so Markov applies.

\textbf{Step 2: Apply Markov}

\[
P (Y \ge a^2) \le \frac{E[Y]}{a^2}
\]

\textbf{Step 3: Substitute back}

\[
P ((X - \mu)^2 \ge a^2) = P (|X - \mu| \ge a)
\]

\textbf{Step 4: Use variance definition}

\[
E[Y] = E[(X - \mu)^2] = \sigma^2
\]

\textbf{Final Result}

\[
P (|X - \mu| \ge a) \le \frac{\sigma^2}{a^2}
\]

\subsection*{2.4 Example}

If:

\[
\mu = 100, \quad \sigma^2 = 25
\]

Then:

\[
P (|X - 100| \ge 10) \le \frac{25}{100} = 0.25
\]

\subsection*{2.5 Insight}

Uses both mean and spread.\\
Stronger than Markov.\\
Still distribution-independent.

\section*{3. Chernoff Bound (Poisson Tail)}

\subsection*{3.1 Motivation}

Markov and Chebyshev give loose bounds.

Need exponentially decreasing bounds for sums of random variables.

\subsection*{3.2 Idea}

Instead of applying Markov to $X$, apply it to:

\[
e^{tX}
\]

This amplifies large deviations.

\subsection*{3.3 Markov on Exponential}

Apply Markov:

\[
P (X \ge a) = P (e^{tX} \ge e^{ta})
\]

\[
P (e^{tX} \ge e^{ta}) \le \frac{E[e^{tX}]}{e^{ta}}
\]

\subsection*{3.4 Resulting Bound}

\[
P (X \ge a) \le \frac{E[e^{tX}]}{e^{ta}}
\]

Choose $t$ to minimize the bound.

\subsection*{3.5 Interpretation}

Uses moment generating function (MGF).\\
Produces exponentially tight bounds.\\
Especially useful for:

Sum of independent random variables\\
Poisson-type tails

\subsection*{3.6 Key Insight}

Transforming the variable makes rare events more visible.

Leads to much stronger guarantees than Chebyshev.

\section*{4. Jensen’s Inequality}

\subsection*{4.1 Convex Functions}

A function $f$ is convex if:

\[
f(\lambda x + (1 - \lambda)y) \le \lambda f(x) + (1 - \lambda)f(y)
\]

for $0 \le \lambda \le 1$.

\subsection*{4.2 Definition of Jensen’s Inequality}

For a convex function $f$:

\[
f(E[X]) \le E[f(X)]
\]

\subsection*{4.3 Proof Idea (Step-by-Step Reasoning)}

\textbf{Step 1: Use convexity}

Convexity says function lies below line segments.

\textbf{Step 2: Represent expectation as weighted sum}

\[
E[X] = \sum_i p_i x_i
\]

\textbf{Step 3: Apply convexity repeatedly}

\[
f\left(\sum_i p_i x_i\right) \le \sum_i p_i f(x_i)
\]

\textbf{Step 4: Recognize expectation}

\[
f(E[X]) \le E[f(X)]
\]

\subsection*{4.4 Example Insight}

If $f(x) = x^2$ (convex):

\[
(E[X])^2 \le E[X^2]
\]

\subsection*{4.5 Interpretation}

Applying function after expectation gives smaller value than:

Taking expectation after applying function

\section*{5. Case Study: Worst-Case Server Demand}

\subsection*{5.1 Problem Context}

Predict maximum load on a server.

Exact distribution unknown.

\subsection*{5.2 Approach Using Inequalities}

\textbf{Step 1: Use expectation}

Apply Markov for rough bound.

\textbf{Step 2: Use variance}

Apply Chebyshev for tighter bound.

\textbf{Step 3: Use exponential bounds}

Apply Chernoff for strong guarantees.

\subsection*{5.3 Key Insight}

These inequalities allow:

Bounding rare events\\
Designing systems with guarantees\\
Understanding worst-case scenarios

\section*{6. Final Summary (Exam Revision)}

\begin{tabular}{|c|c|c|}
\hline
Inequality & Uses & Strength \\
\hline
Markov & Expectation & Weak \\
\hline
Chebyshev & Mean + Variance & Moderate \\
\hline
Chernoff & Exponential moments & Strong \\
\hline
Jensen & Convexity & Structural \\
\hline
\end{tabular}

\subsection*{Final Takeaway}

All inequalities aim to bound:

\[
P (\text{large deviation})
\]

Increasing strength:

Markov $\rightarrow$ Chebyshev $\rightarrow$ Chernoff

Jensen provides a fundamental tool for expectation transformations.

\end{document}